%=====================================================================
% Nomenclature: the symbols the manuscript uses, in one place, with the value
% each takes in the benchmark. Two copies exist on purpose:
%   * this file, for the journal's separate-items upload;
%   * paper/tables/tab_nomenclature.tex, generated from the problem's own
%     parameters so the values cannot drift from the code.
%=====================================================================
\documentclass[11pt]{article}
\usepackage[a4paper,margin=2.5cm]{geometry}
\usepackage{amsmath,amssymb}
\usepackage{booktabs}
\pagestyle{empty}

\begin{document}

\begin{center}
{\large\bfseries Nomenclature}
\end{center}

\vspace{0.5em}
\noindent\textbf{State and control.} $x \in \mathbb{R}^d$ is the state (inventory, and the signal
when one is modelled); $u \in \mathbb{R}^m$ is the control (trading rate); $d$ is the state
dimension and $m$ the control dimension, with $m = p$ and $d = p$ in the inventory-only benchmark;
$p$ is the number of assets.

\medskip
\noindent\textbf{Problem data.} $T$ is the horizon; $Q = \lambda\sigma^2 I$ the inventory-risk
penalty; $R = \eta I$ the temporary market impact; $Q_T$ the terminal inventory penalty;
$\Sigma = \sigma I$ the diffusion; $A$, $B$ the dynamics; $\lambda$ risk aversion (a preference, not
measurable); $\sigma$ relative volatility; $\eta$ the impact coefficient, identified with half the
Roll effective spread; $k = \sqrt{Q/\eta}$ the liquidation rate.

\medskip
\noindent\textbf{Solution.} $V(x,t) = -\tfrac12 x^\top P(t)x + c(t)$ is the value function;
$P(t)$ solves the matrix Riccati equation with $P(T) = Q_T$; $u^\star = -R^{-1}B^\top P(t)x$ is the
optimal feedback; $J(u)$ is the objective under control $u$.

\medskip
\noindent\textbf{Discretisation.} $h$ is the mesh spacing, $N$ the number of grid points per axis,
$N^\star$ the smallest $N$ reaching the stated accuracy target; $R_h$ is the reconstruction that
turns nodal values into a continuous function, of polynomial degree $k$; $\nabla_h V$ is the gradient
recovered from it; $\varepsilon$ is the finite-difference probe step, taken as
$\varepsilon = c\,h^{\alpha}$ when the scaling is being tested.

\medskip
\noindent\textbf{Errors.} All errors are relative $L^2$ norms,
$\|a - a^\star\|_2 / \|a^\star\|_2$; every convergence order is the least-squares slope of
$\log(\text{error})$ against $\log h$; $R^2$ is reported as a secondary, dimensionless accuracy
number and is never used to fit an order.

\medskip
\noindent\textbf{Measured configurations.} Default coefficients $\eta = \lambda\sigma^2 = 0.01$,
$Q_T = 0.1$, $\sigma = 0.01$, $T = 1$ (giving $k = 1$); accuracy-matched grid resolution
$N^\star = 21$ at every dimension for the quadratic benchmark; CPU for grid timings, one GPU for the
neural ones.

\end{document}
